\documentclass[11pt]{article}

\usepackage[a4paper,margin=2.2cm]{geometry}
\usepackage{amsmath,amssymb}
\usepackage{bm}
\usepackage{siunitx}
\usepackage{xcolor}
\usepackage[colorlinks=true,linkcolor=blue]{hyperref}

\setcounter{MaxMatrixCols}{20}
\newcommand{\x}{\bm{x}}
\newcommand{\uu}{\bm{u}}
\newcommand{\vT}{^{\!\top}}
% compact trig shorthand
\newcommand{\cf}{c_\phi}\newcommand{\sfi}{s_\phi}
\newcommand{\cth}{c_\theta}\newcommand{\sth}{s_\theta}\newcommand{\tth}{t_\theta}
\newcommand{\cps}{c_\psi}\newcommand{\sps}{s_\psi}

\title{\textbf{Apollo Lunar Module --- One-Shot Motion-Planning OCP}\\[2pt]
\large Optimal Control Problem Formulation}
\author{}
\date{}

\begin{document}
\maketitle
\vspace{-2.5em}

\section*{State vector (17)}

\[
\x =
\big[\;
\underbrace{x_E,\,y_E,\,z_E}_{\text{position }\bm r_E},\;
\underbrace{u,\,v,\,w}_{\text{body vel. }\bm v_b},\;
\underbrace{\phi,\,\theta,\,\psi}_{\text{Euler }\bm\Theta},\;
\underbrace{p,\,q,\,r}_{\text{body rates }\bm\omega},\;
\underbrace{T,\;\delta_p,\dot\delta_p,\;\delta_y,\dot\delta_y}_{\text{actuator states}}
\;\big]\vT
\]

Frames: $E$ is the (flat-Moon) surface frame, $B$ the body frame ($x$ fwd,
$y$ right, $z$ down); $z_E\le 0$ is above ground.

\section*{Input vector (19)}

\[
\uu =
\big[\;
T_c,\;\delta_{p,c},\;\delta_{y,c},\;
f_1,\,f_2,\,\dots,\,f_{16}
\;\big]\vT
\]

where $T_c,\delta_{p,c},\delta_{y,c}$ are the commanded thrust and pitch/yaw
gimbal angles, and $f_1,\dots,f_{16}$ are the individual RCS thruster forces.

\section*{State equations \ $\dot\x = f(\x,\uu)$}

With the shorthand $c_\ast=\cos\ast$, $s_\ast=\sin\ast$, $t_\ast=\tan\ast$, and
$m,\,J=\operatorname{diag}(I_{xx},I_{yy},I_{zz}),\,g$ the mass, inertia and lunar
gravity, the continuous dynamics are

\begin{align}
\dot{\bm r}_E &= C_{E/B}(\bm\Theta)\,\bm v_b
   && \text{(navigation)}\\[4pt]
\dot{\bm v}_b &= -\,\bm\omega\times\bm v_b \;+\; \bm g_b
   \;+\; \tfrac{1}{m}\big(\bm F_T+\bm F_{\mathrm{rcs}}\big)
   && \text{(translation, body frame)}\\[4pt]
\dot{\bm\Theta} &= E(\bm\Theta)\,\bm\omega
   && \text{(attitude kinematics)}\\[4pt]
\dot{\bm\omega} &= J^{-1}\big(\bm M-\bm\omega\times J\bm\omega\big)
   && \text{(Euler's equations)}
\end{align}

closed by the actuator states
\[
\dot T=\frac{T_c-T}{\tau_T},
\qquad
\ddot\delta_p=\omega_n^2(\delta_{p,c}-\delta_p)-2\zeta\omega_n\dot\delta_p,
\qquad
\ddot\delta_y=\omega_n^2(\delta_{y,c}-\delta_y)-2\zeta\omega_n\dot\delta_y .
\]

\noindent\textbf{Rotation matrix} (body $\to$ surface, 3--2--1 Euler):
\[
C_{E/B}=
\begin{bmatrix}
\cth\cps & \sfi\sth\cps-\cf\sps & \cf\sth\cps+\sfi\sps\\[2pt]
\cth\sps & \sfi\sth\sps+\cf\cps & \cf\sth\sps-\sfi\cps\\[2pt]
-\sth    & \sfi\cth            & \cf\cth
\end{bmatrix},
\qquad
E=
\begin{bmatrix}
1 & \sfi\tth & \cf\tth\\[2pt]
0 & \cf      & -\sfi\\[2pt]
0 & \sfi/\cth & \cf/\cth
\end{bmatrix}.
\]

\noindent\textbf{Forces and moments.} The gimballed engine thrust, gravity, and
the coupled Coriolis term expand (component-wise) as
\[
\bm F_T=
\begin{bmatrix} T\,s_{\delta_p}\\[2pt] -T\,c_{\delta_p}s_{\delta_y}\\[2pt] -T\,c_{\delta_p}c_{\delta_y}\end{bmatrix}\!,
\quad
\bm g_b= g\begin{bmatrix}-\sth\\ \sfi\cth\\ \cf\cth\end{bmatrix}\!,
\quad
\bm\omega\times\bm v_b=
\begin{bmatrix} q w-r v\\ r u-p w\\ p v-q u\end{bmatrix}\!,
\quad
\bm\omega\times J\bm\omega=
\begin{bmatrix}(I_{zz}-I_{yy})\,qr\\ (I_{xx}-I_{zz})\,rp\\ (I_{yy}-I_{xx})\,pq\end{bmatrix}\!.
\]
The RCS force $\bm F_{\mathrm{rcs}}$ and total moment $\bm M$ are
\[
\begin{bmatrix}\bm F_{\mathrm{rcs}}\\[2pt]\bm M_{\mathrm{rcs}}\end{bmatrix}
= B\,\bm f,
\qquad
\bm M=\underbrace{\begin{bmatrix}-d_z\,T_y\\ d_z\,T_x\\ 0\end{bmatrix}}_{\text{gimbal offset }d_z=\SI{2.5}{m}}
+\;\bm M_{\mathrm{rcs}},
\]
where $T_x,T_y$ are the first two components of $\bm F_T$, and $B$ is the fixed
$6\times16$ RCS allocation matrix
\[
B=\big[\; \bm d_i \;;\; \bm r_i\times\bm d_i \;\big]_{i=1}^{16},
\qquad
\bm r_i=\text{thruster position},\quad \bm d_i=\text{unit fire direction}.
\]

\section*{Discretisation}

Direct multiple shooting on a uniform grid of $N=80$ intervals,
$\Delta t=\SI{1}{s}$ (horizon $\SI{80}{s}$), with a 4th-order Runge--Kutta
integrator $\Phi_{\mathrm{RK4}}$:
\[
\x_0 = \x_{\mathrm{init}},
\qquad
\x_{k+1} = \Phi_{\mathrm{RK4}}\!\big(\x_k,\uu_k;\Delta t\big),
\quad k=0,\dots,N-1 .
\]

\section*{Cost function}

\[
\begin{aligned}
J =\;&
\sum_{k=0}^{N-1}
\Big[
(\x_k-\x_{\mathrm{ref}})\vT Q_s (\x_k-\x_{\mathrm{ref}})
+(\uu_k-\uu_{\mathrm{ref}})\vT R\, (\uu_k-\uu_{\mathrm{ref}})
\Big]\\
&+\sum_{k=1}^{N-1}
\Delta\uu_k\vT R_\Delta\,\Delta\uu_k
\;+\;(\x_N-\x_{\mathrm{ref}})\vT Q_f (\x_N-\x_{\mathrm{ref}})
\end{aligned}
\]

with $\Delta\uu_k=\uu_k-\uu_{k-1}$. The middle sum is a \emph{control-rate}
(move-suppression) penalty: unlike $R$, which penalises command \emph{magnitude},
$R_\Delta$ penalises how fast the commands \emph{change}, damping gimbal-command
oscillation. The references are $\x_{\mathrm{ref}}=\bm 0$ except the terminal
altitude (see cut-off, below), and
\[
\uu_{\mathrm{ref}}=\big[\,T_{\mathrm{hover}},\,0,\,0,\,0,\dots,0\,\big]\vT,
\qquad T_{\mathrm{hover}}=m\,g .
\]
Diagonal weights (element-wise; actuator states carry zero weight), grouped
as $[\,\text{pos}\mid\text{vel}\mid\text{att}\mid\text{rate}\mid\text{act}\,]$:
\begingroup\small
\begin{align*}
Q_s &= \operatorname{diag}\!\big(\,20,20,30 \,\big|\, 60,60,60 \,\big|\, 20,20,1 \,\big|\, 30,30,30 \,\big|\, 0,0,0,0,0\,\big)\\[3pt]
Q_f &= \operatorname{diag}\!\big(\,5000,5000,8000 \,\big|\, 4000,4000,4000 \,\big|\, 6000,6000,400 \,\big|\, 10^4,10^4,10^4 \,\big|\, 0,0,0,0,0\,\big)\\[3pt]
R   &= \operatorname{diag}\!\big(\,10^{-7};\; 8,8;\; \underbrace{5\!\cdot\!10^{-4},\dots}_{16}\,\big),
\qquad
R_\Delta = \operatorname{diag}\!\big(\,10^{-6};\; 4000,4000;\; \underbrace{10^{-3},\dots}_{16}\,\big)
\end{align*}
\endgroup

\section*{Path constraints}

\begingroup\small
\[
\begin{aligned}
\text{state } (\forall k=0,\dots,N):\quad
& z_E \le 0, &&
-V_{\max}\le u,v,w\le V_{\max}, & V_{\max}&=\SI{60}{m/s}\\
& -\Theta_{\max}\le \phi,\theta,\psi\le \Theta_{\max}, && \Theta_{\max}=45^\circ, &&\\
& -\Omega_{\max}\le p,q,r\le \Omega_{\max}, && \Omega_{\max}=10^\circ/\text{s}, &&\\
& T_{\min}\le T\le T_{\max}, && -\delta_{\max}\le \delta_p,\delta_y\le \delta_{\max}, & \delta_{\max}&=6^\circ\\[4pt]
\text{input } (\forall k=0,\dots,N-1):\quad
& T_{\min}\le T_c\le T_{\max}, && -\delta_{\max}\le \delta_{p,c},\delta_{y,c}\le \delta_{\max}, &&\\
& 0\le f_i\le F_{\mathrm{rcs}},\ i=1,\dots,16, && F_{\mathrm{rcs}}=\SI{445}{N}, &&
\end{aligned}
\]
\endgroup
\noindent with $[T_{\min},T_{\max}]=[4560,\,45040]\ \si{N}$.

\section*{Two-phase landing with engine cut-off}

The powered OCP above targets a low \emph{contact} altitude
$z_{E}=-h_{\mathrm{contact}}$ ($h_{\mathrm{contact}}=\SI{1}{m}$) over the pad,
rather than the ground. At contact the DPS is cut and the last stretch to the
surface is an \emph{analytic}, engine-off ballistic settle: from the powered
terminal state the thrust is hard-cut ($T:=0$, gimbals zeroed, RCS off) and
$\dot\x=f(\x,\bm 0)$ is integrated with fine steps until $z_E=0$. For a near-rest
contact the touchdown speed is
\[
v_{\mathrm{td}} \;=\; \sqrt{\,w_{\mathrm{c}}^2 + 2\,g\,h_{\mathrm{contact}}\,},
\]
$w_{\mathrm{c}}$ being the vertical descent rate at cut-off. This mirrors the real
DPS shutdown at contact and keeps the OCP well-posed near $z=0$ (no thrust gating
or hover pathology).

\vspace{0.6em}
\noindent\textbf{Remarks.}
\begin{itemize}\itemsep2pt
  \item Thrust and gimbal limits appear \emph{twice} --- as state bounds on the
        actual $T,\delta_p,\delta_y$ and as input bounds on the commands ---
        because the actuators are dynamic states driven by their commands.
  \item The attainable RCS wrench set is the image of the per-thruster box
        $0\le f_i\le F_{\mathrm{rcs}}$ under $B$ (a coupled zonotope), not an
        axis-aligned box.
  \item There is \emph{no} terminal equality constraint; the final state is driven
        only softly by $Q_f$ toward the contact condition.
\end{itemize}

\end{document}
